\subsection{Public boundaries and sequential composition}
\label{sec:composition}

A complete conversation exports a typed record. That record can be the input
to a later conversation. The boundary is explicit: the next endpoints may use
the completed trace and returned values because the interface publishes them.
Only information declared by the return interface crosses that boundary.

\paragraph{The terminal information boundary.}
Each role chooses its terminal return from the preceding trace before seeing
the other roles' returns. Publishing their tuple therefore adds a simultaneous
decision stage when the result types admit choices. It is an additional
imperfect-information feature of the conversation language, even though the
preceding choices were public. Serializing those returns as ordinary public
choices would enlarge later roles' information. When all returns are unit,
the stage carries no decision, as assumed in \cref{sec:rational}.
Throughout this section the boundary has this
explicit simultaneous interpretation; sequencing is composition in the
conversation language, rather than a cut rule for scripts.

A conversation is a finite script, a skip, or a boundary sequence of
conversations. In this section $G$ and $H$ range over conversations. Extend
$\Local{i}{-}$, $\Exec$, and $\Run$ from scripts by the following clauses.

\begin{definition}[Boundary sequencing]
If $G$ is a conversation in $\Gamma$ and
$H$ is a conversation in $\Gamma,z:\Exec(G)$, write $G\mathbin{\triangleright_z}H$
for their boundary sequence. A plan for role $i$ consists of
\begin{equation}
 \Local{i}{G\mathbin{\triangleright_z}H}
   :=\Local{i}{G}\times\prod_{e:\Exec(G)}\Local{i}{H[e/z]}.
 \label{dep:eq:17}
\end{equation}
Write its two components as $(p_i,q_i)$.
Execute $G$ with $\vecp$, obtaining $e$, then execute $H[e/z]$ with
$(q_i(e))_i$. Retain both execution records:
\begin{equation}
 \Exec(G\mathbin{\triangleright_z}H)
       :=\sum_{e:\Exec(G)}\Exec(H[e/z]).
 \label{dep:eq:18}
\end{equation}
\end{definition}

For skip, set every local plan type and its execution type to $\Unit$, and
return $\star$. A component supplies its next-stage plan as a function of the
record that its interface will actually receive.

\begin{proposition}[Sequential execution]
\label{prop:sequence}
The composite is well typed, and its execution satisfies
\begin{equation}
 \Run_{G\triangleright H}(\vecp,\vec q)
   =\bigl(e,\Run_{H[e/z]}((q_i(e))_i)\bigr),
 \qquad e=\Run_G(\vecp).
 \label{dep:eq:19}
\end{equation}
\end{proposition}
\begin{proof}
Induct on conversation formation, using \cref{thm:interaction} for a script
and unit introduction for skip. In the sequence case the first induction
hypothesis gives $e:\Exec(G)$. Dependent application gives
$q_i(e):\Local{i}{H[e/z]}$. The second induction hypothesis produces the
second component of \eqref{dep:eq:18}. Equation \eqref{dep:eq:19} is the definition of the composite
execution.
\end{proof}

For the next statement, compare interface functions pointwise and dependent
pairs componentwise. This extensional interface equality is metatheoretic;
it adds no conversion rule to $\D$. An internal proof that the maps below are
inverse would require the corresponding dependent extensionality principles
in the chosen data calculus.

\begin{lemma}[Associativity of public boundary composition]
\label{thm:associativity}
For $G$, $H(e)$, and $J(e,f)$, the two bracketings of sequential composition
have canonically isomorphic endpoint-plan and execution types at extensional
interface equality. Under those isomorphisms their executions agree.
\end{lemma}
\begin{proof}
For each role, the left-associated plan has components
\[
 \left(p_i,\ q_i:\prod_e\Local{i}{H(e)},\
                  r_i:\prod_{(e,f)}\Local{i}{J(e,f)}\right).
\]
Reassociate it as
$\bigl(p_i,\lambda e.(q_i(e),\lambda f.r_i(e,f))\bigr)$.
The inverse extracts the two functions by projection and currying. The maps
are pointwise inverse by the dependent pair and function laws. The execution types are
related by
\[
 ((e,f),g)\longleftrightarrow(e,(f,g)).
\]
In either bracketing, \eqref{dep:eq:19} first computes $e$ from $\vecp$, then $f$ from
$\vec q(e)$, then $g$ from $\vec r(e,f)$. The three computations and their
dependent arguments are identical.
\end{proof}

A skip conversation is a unit for this composition. The left and
right unit isomorphisms follow from $\Unit\times A\cong A$ and the
corresponding dependent sum laws. Its right unit discards only the extra unit
record, not the execution of the preceding conversation.

\paragraph{What the boundary retains.}
Equations \eqref{dep:eq:18}--\eqref{dep:eq:19} retain intermediate choices and results. Objectives may
inspect the whole composite record. Quotienting the first stage to a Boolean
success value would generally give the second stage a different interface.
Similarly, removing the boundary and concatenating interactions is justified
only by a further translation that preserves when the returned information
becomes available. The public boundary is part of this composition operation.

\begin{proposition}[Definability over the data calculus]
\label{prop:definability}
The endpoint interpretations, complete-execution types, and finite public
execution and sequencing operations are definable in $\D$. Adding their
interface annotations and defining equations is a conservative definitional
extension for judgments in the original data language.
\end{proposition}
\begin{proof}
Expand \eqref{dep:eq:4}, \eqref{dep:eq:9}--\eqref{dep:eq:13}, and \eqref{dep:eq:17}--\eqref{dep:eq:19} recursively on each finite script or
conversation. The result uses only data types, dependent sums and products,
pair elimination, and application in $\D$. Each endpoint constructor expands
by \eqref{dep:eq:5}--\eqref{dep:eq:7}. The translation is the identity on the original data syntax and
types, so a derivation of an original data judgment translates to a derivation
in $\D$ of that same judgment.
\end{proof}

This definability gives the public core a direct implementation route through
an existing dependent kernel. The global annotations expose the ownership and
binding discipline, while the interpretation explains their computational
content. Extending composition to arbitrary open channel interfaces requires
additional syntax and proof rules, considered in \cref{sec:development}.
